\documentclass[withoutpreface]{cumcmthesis}

\usepackage{amsmath,amssymb}
\usepackage{booktabs}
\usepackage{graphicx}
\usepackage{float}
\usepackage{tabularx}
\usepackage{threeparttable}
\usepackage{url}
\usepackage{enumitem}
\usepackage{array}
\usepackage{multirow}
\usepackage{longtable}
\usepackage[table]{xcolor}
\usepackage{hyperref}
\graphicspath{{figures_generated_v5/}}
\hypersetup{colorlinks=true,linkcolor=black,citecolor=black,urlcolor=blue}
\definecolor{TableHeader}{HTML}{DCEAF7}
\definecolor{TableAlt}{HTML}{EEF5FB}
\definecolor{TableAccent}{HTML}{4C78A8}

\title{基于稳健趋势外推、策略级推断与局部鲁棒优化的锂离子电池快充研究}
\tihao{A}
\baominghao{}
\schoolname{}
\membera{}
\memberb{}
\memberc{}
\supervisor{}
\yearinput{2026}
\monthinput{08}
\dayinput{16}

\newcommand{\SOH}{\mathrm{SOH}}
\newcommand{\EOL}{\mathrm{EOL}}
\newcommand{\RMSE}{\mathrm{RMSE}}
\newcommand{\MAE}{\mathrm{MAE}}
\newcommand{\TS}{\mathrm{TS}}

\begin{document}
\maketitle

\begin{abstract}
针对仅给出前200圈实验记录、且部分待预测电池仅公开前150圈数据的锂离子电池快充策略评价问题，本文建立“稳健预处理--早期退化表征--策略级统计推断--样本外预测”的统一框架。全部结论仅使用题目提供的两份CSV数据，不引入被题目删去的真实寿命标签。

首先，以题目给出的初始容量重算健康状态（SOH），用Hampel/MAD规则识别局部异常，仅修复具有充分局部证据的孤立尖峰，再采用窗口11、二阶Savitzky--Golay滤波。将五个稳定段起点与线性、单调幂律、中心化单调二次组成15个候选，在第151--200圈验证，并用外层留一电池评价完整选型流程。一标准误近优集合内再依次比较无效EOL、截断更新、参数边界和复杂度，最终选择$n_0=31$的中心化单调二次；40个外层折中38折选择该结构，另2折选择$n_0=21$中心化二次，外层$\RMSE=4.73\times10^{-4}$。九策略的直接观测稳定退化速率存在显著总体差异（100000次标签置换$p=0.000200$，$\eta_H^2=0.413$），中心化二次EOL的总体差异也达到5\%水平（$p=0.01522$），但Dunn--Holm后没有EOL策略对显著。4.8C--80\%--4.8C新结构的EOL中位数最高，且稳定速率、$S_{150}$和AUC均同向最优，故作为典型长寿命策略；同参数原结构的EOL最短且早期指标较差，作为典型短寿命策略。

其次，从两阶段恒流充电过程推导0--80\%理论时间、平均倍率$A$及低/中高SOC倍率暴露差$D$。数据集3的六种新结构策略中，以稳定退化速率为主响应，720种精确全排列得到$p=13/720=0.0181$，留一策略六折的$A,D$系数均保持正向；随后才报告描述拟合度$R^2=0.998$、调整$R^2=0.997$。这表明较高总体倍率及更偏向中高SOC的倍率暴露与更快退化相关。EOL辅助回归虽有$p=33/720=0.0458$，但Bootstrap系数区间均跨0、留一预测Spearman仅0.371，故仍只作方向辅助，不作因果解释。重参数化后$A,D$的VIF均为1.84；真正的限制是仅有六个不同参数位置、残差自由度仅3。

再次，按题设训练/测试标识，采用“外层留一电池、内层完成全部候选选择”的嵌套验证。40个外层折均在one-SE近优集合中选择预先指定的稳健Theil--Sen基线；预测151--200圈的外层$\MAE=3.43\times10^{-4}$、$\RMSE=5.29\times10^{-4}$、$R^2=0.9970$，相对吻合度为99.965\%。Trend--PCA--Ridge最佳RMSE为$5.83\times10^{-4}$，未超过简单基线。9块测试电池的80\%寿命继承问题一的中心化二次结构，并在对数尺度以75\%个体、25\%同策略训练中心进行部分池化；40折均选择该权重，使EOL截断更新中位差由11.77\%降至6.72\%、均值由30.66\%降至17.99\%。未来斜率辅助头按one-SE简约原则选择全体均值$\Delta k$，只解释总体加速，不参与EOL门控。完整寿命IR、温度和充电时间均值的吻合度为98.74\%、98.25\%和95.65\%，但这些都不是真实EOL准确率。

最后，以49块电池前150圈的观测充电时间和稳定退化速率比较九种真实策略，得到三点经验Pareto前沿。连续优化仅限于数据集3六种新结构策略的原始参数凸包与$(A,D)$特征凸包交集，并施加标准化近邻距离和$T_{0\rightarrow80}\le10.0132$ min约束。在倍率0.1C、SOC 1\%的597个可行网格中，以策略内Bootstrap退化预测的90\%分位数为目标；独立5000次验证Bootstrap表明局部候选5.1C--45\%--4.5C优于最佳已有策略的概率仅47.1\%，配对差90\%分位数为$3.24\times10^{-6}>0$，故不强行推荐新参数。最终推荐已有实验策略5.3C--54\%--4.0C新结构；其在40块完整电池第151--200圈真实退化速率的策略级排序中位列第2，与问题三短期结果没有明显矛盾。该推荐限于当前电芯、实验类别和约10分钟充至80\%的局部设计域。

\keywords{锂离子电池\quad 快充策略\quad Theil--Sen估计\quad 留一电池验证\quad Pareto前沿\quad 鲁棒优化}
\end{abstract}

\section{问题重述与研究目标}

题目给出49块磷酸铁锂电池在九种快充策略下的循环记录。单圈变量包括放电容量、SOH、充电时间、内阻及平均温度；电池级信息包括两阶段充电倍率、切换SOC和数据集编号。研究目标为：

\begin{enumerate}[label=\arabic*)]
  \item 清洗并刻画不同快充策略的早期退化特征，估计达到$\SOH=0.8$时的循环寿命，比较典型长、短寿命策略；
  \item 研究第一阶段倍率$C_1$、切换SOC $q$和第二阶段倍率$C_2$对寿命及衰减速度的影响；
  \item 用训练电池的完整200圈记录预测9块测试电池第151--200圈SOH，并给出寿命估计；
  \item 在观测充电时间与退化速率的经验Pareto比较基础上，于已有实验邻域内完成约10分钟充电约束下的鲁棒策略优化与推荐。
\end{enumerate}

已有研究表明早期循环数据可用于电池寿命预测\cite{severson2019}，但本题观测窗口内绝大多数电池的SOH仍接近1，$0.8$阈值远离数据区间。因此，本文严格区分“151--200圈短期预测”和“80\%阈值远期外推”：前者可由留出区间直接检验，后者只能作为同一数据口径下的\textbf{基准估计循环寿命}，不能视作真实EOL标签。

\section{数据理解、基本假设与符号}

\subsection{数据划分}

\texttt{cycle\_train.csv}共有9350行，覆盖49块电池。训练组40块电池具有1--200圈记录；测试组9块电池仅以1--150圈作为可见输入。所有预处理阈值、特征和模型选择均先在训练组内确定；测试组第151--200圈不参与任何模型选择，以避免信息泄漏。

\subsection{基本假设}

\begin{enumerate}[label=H\arabic*：]
  \item 题目记录的容量、内阻、温度和时间量纲一致；\texttt{battery\_summary.csv}中的初始容量可作为各电池SOH归一化基准。
  \item 单个孤立突跳主要来自采集或记录误差；连续变化视为可能的真实状态演化，不做批量替换。
  \item 在第1--200圈的局部窗口内，近期退化趋势可用于预测随后50圈；超过该窗口的阈值寿命仅作模型外推。
  \item 参数效应分析以不同“策略参数位置”为分析单位；同一策略内多块电池是重复样本，不能伪装成新增参数组合。
  \item 根据理论时间与记录量级的一致性，将\texttt{chargetime}解释为逐循环充至约80\% SOC的实际快充时间；该判断是数据口径验证，不等同字段定义的数学证明。
\end{enumerate}

\begin{table}[H]
\centering
\rowcolors{2}{TableAlt}{white}
\begin{tabularx}{\textwidth}{lX}
\toprule
\rowcolor{TableHeader}
符号 & 含义 \\
\midrule
$Q_{i,0}$ & \texttt{battery\_summary.csv}中电池$i$的初始容量 \\
$Q_{i,n}$ & 电池$i$第$n$圈放电容量 \\
$S_{i,n}$ & 电池$i$第$n$圈清洗平滑后的SOH \\
$C_1,q,C_2$ & 第一阶段倍率、切换SOC比例和第二阶段倍率 \\
$L_i$ & 电池$i$达到$\SOH=0.8$的基准估计循环寿命 \\
$n_0$ & 长期稳定段拟合起点，本轮正式值为31 \\
$A_i,B_i,C_i$ & 电池$i$中心化单调二次的水平、线性退化与加速退化参数 \\
$w$ & Q3对数EOL部分池化的个体权重，本轮正式值为0.75 \\
$r_i^{\mathrm{stable}}$ & 稳定段Theil--Sen衰减速率，即斜率的相反数 \\
$R_g$ & 策略$g$内$r_i^{\mathrm{stable}}$的电池级中位数 \\
$A$ & 0--80\%区间按SOC长度加权的平均倍率 \\
$E_L,E_H$ & 低SOC和中高SOC区间的平均倍率暴露 \\
$D$ & 倍率暴露差，$D=E_H-E_L$ \\
$\ell$ & 第三问可见的早期循环长度 \\
$W_{\ell}$ & 截止到第$\ell$圈时的最近30圈窗口 \\
$h$ & 未来预测步长，$h=1,\ldots,50$ \\
$X_{i,n}$ & IR、平均温度或充电时间中的任一状态通道 \\
$T_i^{\mathrm{obs}}$ & 电池$i$前150圈\texttt{chargetime}中位数 \\
$T_{0\rightarrow80}^{\mathrm{theo}}(\boldsymbol\theta)$ & 两阶段恒流从0充至80\% SOC的理论时间 \\
$\boldsymbol\theta=(C_1,q,C_2)$ & 问题四的两阶段策略决策向量 \\
$\Omega_{\mathrm{local}}$ & 原始参数域、$(A,D)$特征域及近邻距离共同限定的局部可行域 \\
$\mathcal P_{\mathrm{obs}}$ & 九种真实策略中基于观测时间与退化速率的经验Pareto前沿 \\
$Q_\alpha(\cdot)$ & 样本经验分布的$\alpha$分位数 \\
\bottomrule
\end{tabularx}
\end{table}
\refstepcounter{table}

\section{统一预处理与评价设计}

\subsection{稳健异常检测和SOH重算}

容量SOH定义为
\begin{equation}
S_{i,n}^{\mathrm{raw}}=\frac{Q_{i,n}}{Q_{i,0}}.
\end{equation}
对每块电池、每个变量分别在局部窗口中计算中位数$m_{i,n}$和
\begin{equation}
\mathrm{MAD}_{i,n}=\operatorname{median}_{j\in W_n}|x_{i,j}-m_{i,n}|.
\end{equation}
局部窗口取当前点前后各5圈（完整窗口共11圈）。若
\begin{equation}
|x_{i,n}-m_{i,n}|>
\max\{5\times1.4826\,\mathrm{MAD}_{i,n},\delta_x\},
\end{equation}
则记为异常候选。容量、内阻、温度和充电时间的最小偏差门槛依次取$0.01Q_{i,0}$、局部典型内阻的2\%、$1\,^{\circ}\mathrm C$和$0.5\,\mathrm{min}$。只有同时满足“单点孤立、相邻点一致、替换后恢复局部连续”的候选才以相邻点均值修复；正值充电时间候选仅标记，不因统计异常自动替换。四类变量分别识别1、18、4、17个候选，实际修复1、16、3、0个孤立点，其中容量异常为1号电池第12圈。该规则保证论文、代码和结果使用同一组$k=5$与最小偏差门槛。

修复容量后重新计算SOH，再使用窗口11、二阶Savitzky--Golay滤波\cite{savitzky1964}得到$S_{i,n}$。图\ref{fig:preprocess}说明先修正尖峰再平滑的必要性：若直接对异常SOH平滑，尖峰会污染相邻多个循环。

\begin{figure}[H]
\centering
\includegraphics[width=0.94\textwidth]{fig1_preprocess.pdf}
\refstepcounter{figure}\label{fig:preprocess}
\end{figure}

\subsection{早期退化指标}

在第1--150圈提取四类核心指标，其中稳健斜率采用Theil--Sen估计\cite{sen1968}：
\begin{align}
S_{150,i}&=\operatorname{median}(S_{i,141},\ldots,S_{i,150}),\\
k_i^{\TS}&=\operatorname{median}_{a<b}\frac{S_{i,b}-S_{i,a}}{b-a},\\
\mathrm{AUC}_i&=\frac{1}{149}\sum_{n=1}^{149}\frac{S_{i,n}+S_{i,n+1}}{2},\\
T_i&=\operatorname{median}_{1\le n\le150}t_{i,n}.
\end{align}
同时保留内阻斜率和平均温度作为辅助诊断。$S_{150}$反映当前健康水平，$k^{\TS}$反映稳健衰减速度，AUC反映整个早期窗口的累计健康状态。九种策略的中位轨迹与四分位区间见图\ref{fig:policysoh}。

\begin{figure}[H]
\centering
\includegraphics[width=0.98\textwidth]{fig2_policy_soh.pdf}
\refstepcounter{figure}\label{fig:policysoh}
\end{figure}

\subsection{统一验证原则}

模型误差使用
\begin{equation}
\MAE=\frac1N\sum_{r=1}^N|y_r-\hat y_r|,
\qquad
\RMSE=\sqrt{\frac1N\sum_{r=1}^N(y_r-\hat y_r)^2}.
\end{equation}
第一问先用每块训练电池第1--150圈拟合、第151--200圈作开发阶段结构比较，再以外层LOBO评价完整选型流程；第三问同样采用外层LOBO，并在每个外层训练集内重新选择趋势基线、特征组、PCA维数和Ridge参数。除短期误差外，还检查参数贴边、单电池误差和远期外推敏感性。外推寿命与$S_{150}$、$k^{\TS}$和AUC的秩一致性只作为辅助诊断：由于线性寿命与线性斜率存在结构性联系，该指标不能作为偏向线性模型的硬判据。

\section{问题一：策略比较与基准寿命估计}

问题一先统一完成异常识别、SOH重算与平滑，再从共同观测窗口提取稳健退化指标，并通过“稳定起点--退化结构”联合验证得到基准寿命。完整求解链如图\ref{fig:workflowq1}所示，图中把数据清洗、指标构造、样本外选型与策略比较放在同一证据链中，避免将远期EOL外推与前150圈直接观测混为一谈。

\begin{figure}[H]
\centering
\includegraphics[width=0.98\textwidth]{fig_workflow_q1.pdf}
\refstepcounter{figure}\label{fig:workflowq1}
\end{figure}

\subsection{稳定段、候选结构与寿命公式推导}

首批循环包含形成过程和局部回升，故把稳定段起点也纳入选择。令
\[
n_0\in\{11,21,31,41,51\},\qquad n=n_0,\ldots,150,
\]
并在每个$n_0$下分别拟合
\begin{align}
\text{线性：}\quad &S_i(n)=a_i+b_i n,\quad b_i<0,\\
\text{单调幂律：}\quad &S_i(n)=a_i-b_i n^{c_i},\quad b_i>0,\ c_i>0,\\
\text{中心化单调二次：}\quad
&S_i(n)=A_i-B_i(n-n_0)-C_i(n-n_0)^2,\quad B_i,C_i\ge0.
\end{align}
阈值寿命由$S_i(L_i)=0.8$逐步求得。线性模型有
\begin{align}
0.8&=a_i+b_iL_i,\\
b_iL_i&=0.8-a_i,\\
L_{i,\mathrm{lin}}&=\frac{0.8-a_i}{b_i}.
\end{align}
幂律模型有
\begin{align}
0.8&=a_i-b_iL_i^{c_i},\\
b_iL_i^{c_i}&=a_i-0.8,\\
L_i^{c_i}&=\frac{a_i-0.8}{b_i},\\
L_{i,\mathrm{pow}}&=\left(\frac{a_i-0.8}{b_i}\right)^{1/c_i}.
\end{align}
中心化二次只从候选稳定起点以后施加单调约束，因为
\begin{equation}
\frac{\mathrm dS_i}{\mathrm dn}=-B_i-2C_i(n-n_0)\le0,\qquad n\ge n_0.
\end{equation}
令$x_i=L_i-n_0$，则阈值方程逐步化为
\begin{align}
A_i-B_ix_i-C_ix_i^2&=0.8,\\
C_ix_i^2+B_ix_i-(A_i-0.8)&=0.
\end{align}
当$C_i>0$时取唯一非负根
\begin{equation}
\boxed{L_{i,\mathrm{cq}}=n_0+
\frac{-B_i+\sqrt{B_i^2+4C_i(A_i-0.8)}}{2C_i}},
\end{equation}
当$C_i=0,B_i>0$时退化为$L_i=n_0+(A_i-0.8)/B_i$。该参数化允许稳定段内退化速度逐渐加快，又不强迫第0圈起就满足下降。
只有参数满足单调约束且交点有限时才保留寿命。先定义稳定段Theil--Sen斜率，再取其相反数作为正向退化速度：
\begin{align}
k_i^{\mathrm{stable}}
&=\operatorname{median}_{n_0\le u<v\le150}
\frac{S_{i,v}-S_{i,u}}{v-u},\\
&\boxed{r_i^{\mathrm{stable}}=-k_i^{\mathrm{stable}}},\qquad
r_g^{\mathrm{stable}}=\operatorname{median}_{i\in g}r_i^{\mathrm{stable}}.
\end{align}
故$r_i^{\mathrm{stable}}$越大表示退化越快，越小表示退化越慢；它直接来自前150圈观测SOH，不依赖远期EOL交点，是问题二与问题四的主退化指标。

\subsection{一标准误选择与外层验证}

对候选$m=(n_0,\text{结构})$，先按电池计算第151--200圈MSE，再求电池级均值$\overline e_m$及标准误$\mathrm{SE}_m$。若$m^*=\arg\min_m\overline e_m$，则近优集合为
\begin{equation}
\mathcal A=\left\{m:\overline e_m\le
\overline e_{m^*}+\mathrm{SE}_{m^*}\right\}.
\end{equation}
在$\mathcal A$中按预先冻结的顺序比较：无效EOL数、前150圈与前200圈拟合同一结构所得EOL的中位相对更新差、参数边界数、结构复杂度与$n_0$。因此截断稳定性是近优集合内的核心长期指标，边界仅作次级诊断；该顺序在外层结果揭示前固定。

\begin{table}[H]
\centering
\resizebox{\textwidth}{!}{%
\rowcolors{2}{TableAlt}{white}
\begin{tabular}{rrrrrl}
\toprule
\rowcolor{TableHeader}
$n_0$ & 线性RMSE & 幂律RMSE & 中心二次RMSE & 中心二次EOL更新中位数 & 说明 \\
\midrule
11 & 0.001254 & 0.000704 & 0.000579 & 8.48\% &  \\
21 & 0.001076 & 0.000485 & 0.000506 & 9.48\% &  \\
31 & 0.000954 & \textbf{0.000423} & 0.000454 & 11.77\% & \textbf{one-SE正式选择} \\
41 & 0.000860 & 0.000428 & 0.000447 & 13.14\% &  \\
51 & 0.000787 & 0.000484 & 0.000482 & 20.00\% &  \\
\bottomrule
\end{tabular}
}
\refstepcounter{table}\label{tab:q1models}
\end{table}

开发集最低误差来自$n_0=31$幂律，但$n_0=31$中心化二次同处one-SE近优集合，且EOL截断更新中位数为11.77\%，明显低于同起点幂律的19.74\%，故按冻结规则正式选择$n_0=31$中心化二次。外层每次留出一块训练电池，并仅在其余39块内重复全部选择：38折选择$n_0=31$中心化二次，2折选择$n_0=21$中心化二次，外层$\MAE=0.000285$、$\RMSE=0.000473$。相较旧$n_0=21$幂律的真正嵌套$\RMSE=0.000483$，下降约2.09\%；图\ref{fig:modelselect}同时给出开发阶段候选与外层频次。热图把15个“稳定起点--退化结构”组合放到同一色阶中，可见线性结构在各起点下的误差整体更高；$n_0=31$中心化二次虽不是开发集绝对最小误差点，却位于一标准误近优区域，并在外层40折中有38折被重新选择，说明其优势主要来自“近期预测误差、远期稳定性与复杂度”的综合平衡，而非单次拟合最优。

这里$n_0=31$表示从该位置开始，各电池总体稳健趋势适合由统一下降结构描述；它不表示所有电池在第31圈后逐点严格单调，更不表示某个电化学过程在该圈绝对结束。图\ref{fig:q1fitshape}在稳定退化速率10\%、50\%和90\%分位附近各取一块完整电池，显示三类候选的实际拟合形态。

\begin{figure}[H]
\centering
\includegraphics[width=0.98\textwidth]{fig4_model_selection.pdf}
\refstepcounter{figure}\label{fig:modelselect}
\end{figure}

\begin{figure}[H]
\centering
\includegraphics[width=0.99\textwidth]{fig4b_q1_fit_shapes.pdf}
\refstepcounter{figure}\label{fig:q1fitshape}
\end{figure}

用前150圈和前200圈分别按同一中心化二次结构外推，40块完整电池的EOL相对更新中位数为11.77\%，但均值达到30.66\%，Spearman相关仅0.847。旧$n_0=21$幂律的对应三项为19.77\%、22.95\%和0.920。故新结构改善了典型电池的截断稳定性与近期样本外拟合，却没有全面改善尾部漂移和跨电池排序；少数电池仍可能产生较大远期变化。本文据此只称其为更平衡的内部外推结构，不把交点视为已观测真实寿命。

\subsection{九种策略的比较结果}

49块电池的策略、前150圈充电时间、$S_{150}$、稳定退化速率及基准EOL已逐块汇总于附录表\ref{tab:q1battery}。正文按策略报告中位数：

\begin{table}[H]
\centering
\small
\resizebox{\textwidth}{!}{
\rowcolors{2}{TableAlt}{white}
\begin{tabular}{lrrrrr}
\toprule
\rowcolor{TableHeader}
策略 & $N$ & $S_{150}$ & $10^4r^{\mathrm{stable}}$ & AUC & EOL中位数/圈 \\
\midrule
4.8C--80\%--4.8C（新结构） & 6 & 0.99802 & 0.181 & 0.99902 & 1775 \\
5.6C--19\%--4.6C（新结构） & 8 & 0.99642 & 0.306 & 0.99841 & 1660 \\
5.3C--54\%--4.0C（新结构） & 8 & 0.99717 & 0.217 & 0.99872 & 1593 \\
5.6C--36\%--4.3C（新结构） & 8 & 0.99676 & 0.273 & 0.99867 & 1593 \\
5.0C--67\%--4.0C（新结构） & 7 & 0.99731 & 0.231 & 0.99885 & 1504 \\
3.6C--80\%--3.6C & 3 & 0.99624 & 0.197 & 0.99724 & 1389 \\
80\%--3.6C & 3 & 0.99634 & 0.272 & 0.99817 & 1114 \\
3.7C--31\%--5.9C（新结构） & 3 & 0.97277 & 1.925 & 0.98705 & 856 \\
4.8C--80\%--4.8C & 3 & 0.99021 & 0.849 & 0.99662 & 640 \\
\bottomrule
\end{tabular}}
\refstepcounter{table}\label{tab:q1policy}
\end{table}

\begin{figure}[H]
\centering
\includegraphics[width=0.95\textwidth]{fig5_life_distribution.pdf}
\refstepcounter{figure}\label{fig:lifedist}
\end{figure}

图\ref{fig:lifedist}采用小提琴密度、四分位区间与电池级散点同时展示策略内分布。对数横轴下，4.8C--80\%--4.8C新结构整体位于更长寿命一侧，而同参数原结构集中在较短寿命区域；部分新结构策略存在较长右尾，说明仅比较均值会掩盖小样本策略内的外推离散性，因此正文排序以中位数并结合早期稳定退化指标共同判断。

题目要求提取典型长、短寿命策略。按问题一统一的中心化二次EOL口径，\textbf{正式典型长寿命策略为4.8C--80\%--4.8C新结构}：其策略级EOL中位数最高，且稳定退化速率、$S_{150}$和AUC三项早期指标也均排第1。4.8C--80\%--4.8C原结构的EOL最短，且三项早期指标整体较差，方向一致，故称为典型短寿命策略。两者来自不同数据集/结构类别，因此只能作综合表现对照，不能将差异归因于同一组数值参数。

\begin{table}[H]
\centering
\small
\rowcolors{2}{TableAlt}{white}
\begin{tabular}{lrr}
\toprule
\rowcolor{TableHeader}
指标 & 正式典型长寿命 & 典型短寿命 \\
\midrule
策略 & 4.8C--80\%--4.8C新 & 4.8C--80\%--4.8C \\
样本数 & 6 & 3 \\
充电时间中位数/min & 11.038 & 10.054 \\
$S_{150}$中位数 & 0.99802 & 0.99021 \\
$10^4r^{\mathrm{stable}}$ & 0.181 & 0.849 \\
SOH AUC中位数 & 0.99902 & 0.99662 \\
平均温度中位数/$^\circ$C & 33.12 & 30.07 \\
EOL中位数/圈 & 1775 & 640 \\
早期三指标排序 & 1/1/1 & 8/8/8 \\
\bottomrule
\end{tabular}
\refstepcounter{table}\label{tab:longshortdetail}
\end{table}

图\ref{fig:longshort}显示长寿命策略在第150圈保持更高SOH且稳定退化速率更低，和EOL排序同向；短寿命策略下降更快。两者跨数据集或结构，故问题一只比较综合表现，不把单个参数、温度或时间差解释为因果效应。

\begin{figure}[H]
\centering
\includegraphics[width=0.94\textwidth]{fig6_long_short.pdf}
\refstepcounter{figure}\label{fig:longshort}
\end{figure}

\section{问题二：快充参数影响分析}

问题二在九策略总体差异检验之后，进一步把原始$C_1,q,C_2$重参数化为总体倍率水平$A$与SOC区间暴露差$D$，并依次进行精确置换、留一策略位置验证、分界敏感性和机制通道审计。整体流程见图\ref{fig:workflowq2}；其中稳定退化速率始终作为直接观测主响应，基准EOL仅作辅助方向证据。

\begin{figure}[H]
\centering
\includegraphics[width=0.98\textwidth]{fig_workflow_q2.pdf}
\refstepcounter{figure}\label{fig:workflowq2}
\end{figure}

\subsection{策略总体差异与事后比较}

问题二先以直接观测得到的$r_i^{\mathrm{stable}}$为主响应，以中心化二次基准EOL $L_i$为辅助响应。对九种策略采用Kruskal--Wallis秩检验\cite{kruskal1952}：
\begin{equation}
H=\frac{12}{N(N+1)}\sum_{g=1}^{G}\frac{R_g^2}{n_g}-3(N+1),
\qquad
\eta_H^2=\frac{H-G+1}{N-G}.
\end{equation}
为降低小样本渐近近似误差，保持各组样本数不变，随机打乱策略标签$B=100000$次。Monte Carlo置换$p$值为
\begin{equation}
p_{\mathrm{MC}}
=\frac{1+\#\{H^{(b)}\ge H_{\mathrm{obs}}\}}{B+1}.
\end{equation}
稳定退化速率得到$H=24.504$、渐近$p=0.00189$，共有19次置换不低于观测值，故$p_{\mathrm{MC}}=20/100001=0.000200$，$\eta_H^2=0.413$。总体检验显著后进行Dunn比较\cite{dunn1964}并用Holm法\cite{holm1979}校正：3.6C--80\%--3.6C与3.7C--31\%--5.9C新结构、以及3.7C--31\%--5.9C新结构与4.8C--80\%--4.8C新结构两对显著，调整后$p$分别为0.0435和0.0275。

辅助EOL响应得到$H=17.097$、渐近$p=0.0291$，1521次置换不低于观测值，$p_{\mathrm{MC}}=1522/100001=0.01522$，$\eta_H^2=0.227$。总体检验达到5\%水平，但Dunn--Holm后没有任何EOL策略对显著，说明总体排序分离尚不能稳定定位到单一两两差异。此处的100000次随机置换与后文六个参数位置的$6!=720$种精确全排列是两件不同的检验。

\subsection{两阶段充电时间与倍率暴露推导}

设切换SOC为$q=Q_1/100$。恒流倍率为$C$时，充入$\Delta s$比例容量所需时间为$60\Delta s/C$分钟。因此0--$q$阶段与$q$--0.8阶段相加，得到
\begin{equation}
T_{0\rightarrow80}=60\left(\frac{q}{C_1}+\frac{0.8-q}{C_2}\right)\ \mathrm{min}.
\end{equation}
题面说明完整充电还包含80\%--100\%统一的1C CC--CV阶段，但数据字段的量级表明\texttt{chargetime}并非0--100\%完整时间：仅以理想1C恒流补充20\%容量就至少需要12 min，而记录值主要为10--13 min。反之，3.6C--80\%--3.6C策略的理论值为13.333 min、观测中位数为13.342 min；九策略理论值与观测值之差的中位数为0.051 min。故本文根据理论--实测一致性，将\texttt{chargetime}解释为逐循环实际充至约80\% SOC的快充时间，记为$T_i^{\mathrm{obs}}$。九策略平均绝对偏差仍有0.153 min，且4.8C--80\%--4.8C新结构偏差达到1.038 min，因此该判断用于统一字段口径，不能把理论式当成无误差的实测预测器。原先的“80--100\%时间代理”不再使用；若计算$\Delta T=T_i^{\mathrm{obs}}-T_{0\rightarrow80}^{\mathrm{theo}}$，它只表示实际快充相对理想恒流公式的偏差。

数据集3六种新结构策略的$T_{0\rightarrow80}^{\mathrm{theo}}$位于9.9900--10.0132 min，说明其参数设计近似等时。这使问题二可在近似相同理论时间下重点研究倍率在SOC区间内的分配，而问题四则把10.0132 min作为局部候选的理论时间上限。

定义全区间平均倍率
\begin{equation}
A=\frac{C_1q+C_2(0.8-q)}{0.8}.
\end{equation}
对给定SOC分界$s_0$，定义低SOC和中高SOC平均暴露
\begin{align}
E_L(s_0)&=\frac{C_1\min(q,s_0)+C_2[s_0-\min(q,s_0)]}{s_0},\\
E_H(s_0)&=\frac{C_1\max(q-s_0,0)+C_2[0.8-\max(q,s_0)]}{0.8-s_0},\\
D(s_0)&=E_H(s_0)-E_L(s_0).
\end{align}
主分析取$s_0=0.5$，并在0.4--0.6范围做敏感性分析。

\subsection{相关性、共线性与策略级回归}

同策略电池共享$C_1,q,C_2$，故电池级相关只作描述，不能把重复电池当作新的参数位置。九种策略中共有八种参数完整；主参数模型主动限制在同时满足\texttt{dataset=3}与\texttt{category=NEWSTRUCTURE}的六个策略位置，目的是尽量控制数据集与结构类别混杂，并非“只有六种参数完整”。其他参数完整策略只用于跨策略描述和敏感性对照。以八个参数完整的策略位置为单位，$C_1,q,C_2$与稳定退化速率的Spearman相关分别为0.108、$-0.561$、0.566；与EOL的相关分别为0.639、$-0.171$、$-0.374$。原始参数$C_1,q,C_2$的VIF分别为5.51、1.99、6.13，说明设计存在耦合，故使用$A,D$重参数化。数据集3六种新结构策略中，$A,D$相关约0.676、VIF均为1.84；相关存在但不严重，当前识别能力主要受六个不同参数位置限制。

为控制数据集和电池结构混杂，正式回归只使用数据集3六种新结构策略。把$A,D$策略级标准化为$A^*,D^*$，以策略稳定退化速率中位数为主响应：
\begin{equation}
r^{\mathrm{stable}}=\beta_0+\beta_A A^*+\beta_DD^*+\varepsilon.
\end{equation}
先报告稳健性证据。对六个响应的全部$6!=720$种排列完全枚举，其中13种的$R^2$不低于观测值：
\begin{equation}
p_{\mathrm{exact}}
=\frac{\#\{R_b^2\ge R_{\mathrm{obs}}^2\}}{720}
=\frac{13}{720}=0.0181.
\end{equation}
留一策略位置六折中$\beta_A,\beta_D$均保持正向，留出预测Spearman相关为0.886。随后报告拟合式
\begin{equation}
\boxed{r^{\mathrm{stable}}
=5.223\times10^{-5}
+4.020\times10^{-5}A^*
+2.823\times10^{-5}D^*},
\end{equation}
其$R^2=0.9983$、调整$R^2=0.9972$。策略内分层Bootstrap 95\%区间分别为$[2.90,4.79]\times10^{-5}$和$[1.91,3.50]\times10^{-5}$。Bootstrap只能传播同一策略重复电池的波动，不能创造新的$A,D$位置；由于$n=6$、残差自由度为$6-3=3$，零假设下$E(R^2)\approx2/(6-1)=0.4$，所以$R^2$只作描述，证据顺序应是“精确置换、LOPO方向、Bootstrap、最后才是$R^2$”。

\begin{table}[H]
\centering
\rowcolors{2}{TableAlt}{white}
\begin{tabular}{lrrrr}
\toprule
\rowcolor{TableHeader}
响应 & $\beta_A$ & $\beta_D$ & 精确$p$ & LOPO方向稳定性 \\
\midrule
$r^{\mathrm{stable}}$（主） & $4.02\times10^{-5}$ & $2.82\times10^{-5}$ & $13/720$ & 6/6、6/6为正 \\
$L$（辅助） & $-235.4$ & $-74.3$ & $33/720$ & 6/6、6/6为负 \\
\bottomrule
\end{tabular}
\refstepcounter{table}\label{tab:q2reg}
\end{table}

辅助EOL回归为$L=1496.8-235.4A^*-74.3D^*$，$R^2=0.952$、调整$R^2=0.920$，精确$p=33/720=0.0458$。然而$\beta_A$与$\beta_D$的Bootstrap区间分别为$[-1921.7,534.9]$和$[-1110.5,1171.5]$，均跨0；LOPO虽六折方向均为负，定量预测Spearman仅0.371、RMSE为133圈。因此它只提供方向辅助，不足以证明参数对真实寿命的显著作用。

\subsection{使用观测SOH斜率进行方向一致性检验}

这里$r^{\mathrm{stable}}$与问题一EOL均由前150圈SOH得到，二者方向一致只能说明“观测退化与寿命外推在不同表征下相容”，不能称为统计独立的重复证据。相对内阻与温度来自不同测量通道，可在后文作为更独立但仍为描述性的机制线索。

\subsection{SOC分界敏感性与关联解释}

将$s_0$从0.40变化到0.60，分别以$E_L,E_H$回归$r^{\mathrm{stable}}$。图\ref{fig:socsens}表明两类标准化系数始终为正，且中高SOC系数始终更大：低SOC系数由$8.55\times10^{-5}$降至$5.13\times10^{-5}$，中高SOC系数由$1.42\times10^{-4}$降至$1.02\times10^{-4}$。因此，在这六个近等时策略位置上，较高总体倍率与更快退化相关，把更高倍率暴露安排到中高SOC一侧与更快退化的关联更强；这仍是相关性陈述，不是电化学因果定律。

\begin{figure}[H]
\centering
\includegraphics[width=0.94\textwidth]{fig7_soc_sensitivity.pdf}
\refstepcounter{figure}\label{fig:socsens}
\end{figure}

\subsection{可能机制及证据边界}

为避免用同一SOH曲线的寿命外推反向“证明”机制，本文只把前150圈观测到的相对内阻斜率和平均温度作为不同测量通道的辅助证据。数据集3六个策略位置的Spearman结果见表\ref{tab:mechanism}。

\begin{table}[H]
\centering
\rowcolors{2}{TableAlt}{white}
\begin{tabular}{llrr}
\toprule
\rowcolor{TableHeader}
解释量 & 响应量 & $\rho$ & 未校正$p$值 \\
\midrule
$A$ & 负SOH斜率 & 0.667 & 0.148 \\
$A$ & 相对内阻增长斜率 & 0.928 & 0.0077 \\
$A$ & 平均温度 & 0.841 & 0.0361 \\
$E_H$ & 相对内阻增长斜率 & $-0.143$ & 0.787 \\
$E_H$ & 平均温度 & $-0.086$ & 0.872 \\
平均温度 & 负SOH斜率 & 0.771 & 0.0724 \\
\bottomrule
\end{tabular}
\refstepcounter{table}\label{tab:mechanism}
\end{table}

平均倍率$A$与内阻增长、平均温度均呈正向关联，且温度与更快SOH衰减方向一致。这与较大充电电流可能伴随更强极化、热负荷和阻抗增长的解释相容；既有充电应力实验也观察到电流或截止状态变化会影响容量衰减和内阻增长\cite{gao2017}。但$E_H$与内阻、温度并无同向证据，故不能声称“中高SOC倍率已经通过升温或内阻增长导致寿命下降”。SOC位置效应目前只由稳定退化回归与分界敏感性支持，具体反应通道尚未识别。

此外，快充影响具有电池体系和协议依赖性：在特定LiFePO$_4$多阶段快充条件下，也有研究观察到与标准充电相近的老化结果\cite{ansean2016}。本题只有六个策略位置，上述相关$p$值未作多重比较校正，且缺少逐段电压、电流和电化学诊断，因此机制部分应表述为“与可能机制一致的描述性线索”，而非中介效应或因果结论。

\section{问题三：未来50圈SOH预测与寿命估计}

问题三严格按“可见早期信息--样本外短期预测--远期结构审计”的层次展开：先比较简单趋势基线，再检验多维增强模型是否具有稳定增益，最后才把冻结的中心化二次结构用于EOL外推与同策略部分池化。图\ref{fig:workflowq3}给出无泄漏求解流程，并突出“短期SOH高精度不等于远期EOL高精度”的证据边界。

\begin{figure}[H]
\centering
\includegraphics[width=0.98\textwidth]{fig_workflow_q3.pdf}
\refstepcounter{figure}\label{fig:workflowq3}
\end{figure}

\subsection{预测任务和无泄漏验证}

训练电池集合记为$\mathcal T$（40块），测试电池集合记为$\mathcal U$（9块）。对每块电池仅使用前$\ell$圈构造特征，预测随后50圈；正式任务取$\ell=150$，早期长度分析另取$\ell\in\{50,75,100,125,150\}$。每个$\ell$下均重新构造特征，不能把第$\ell$圈以后的信息带入模型。

\subsection{趋势基线与Trend--PCA--Ridge增强模型}

先定义四种不使用跨电池学习的简单基线。末端稳健状态为
\begin{equation}
\bar S_{i,\ell}=\operatorname{median}
\{S_{i,\ell-9},\ldots,S_{i,\ell}\},
\end{equation}
最近30圈窗口为
\begin{equation}
W_{\ell}=\{\max(1,\ell-29),\ldots,\ell\}.
\end{equation}
在该窗口内，Theil--Sen近期斜率逐步定义为
\begin{equation}
k_{i,\ell}^{\TS}
=\operatorname{median}_{u<v,\ u,v\in W_{\ell}}
\frac{S_{i,v}-S_{i,u}}{v-u}.
\end{equation}
于是未来第$h$步（$h=1,\ldots,50$）的TS基线为
\begin{equation}
\widehat S_{i,\ell+h}^{\TS}
=\bar S_{i,\ell}+h k_{i,\ell}^{\TS}.
\label{eq:tsbaseline}
\end{equation}
OLS基线在同一窗口内使用
\begin{equation}
k_{i,\ell}^{\mathrm{OLS}}
=\frac{\sum_{n\in W_{\ell}}(n-\bar n)(S_{i,n}-\bar S_W)}
{\sum_{n\in W_{\ell}}(n-\bar n)^2},
\qquad
\widehat S_{i,\ell+h}^{\mathrm{OLS}}
=\bar S_{i,\ell}+h k_{i,\ell}^{\mathrm{OLS}}.
\end{equation}
Persistence基线令$\widehat S_{i,\ell+h}=\bar S_{i,\ell}$。

SG导数基线则在每个中心$n$的11圈局部窗口上拟合二阶多项式
\begin{equation}
(\widehat\alpha_0,\widehat\alpha_1,\widehat\alpha_2)
=\arg\min_{\alpha_0,\alpha_1,\alpha_2}
\sum_{j=-5}^{5}\left[S_{i,n+j}-(\alpha_0+\alpha_1j+\alpha_2j^2)\right]^2.
\end{equation}
因为$\mathrm d(\alpha_0+\alpha_1j+\alpha_2j^2)/\mathrm dj|_{j=0}=\alpha_1$，故$\widehat\alpha_1$就是局部一阶导数；取末20圈导数中位数$d_{i,\ell}$后，得到$\widehat S_{i,\ell+h}^{D}=\bar S_{i,\ell}+hd_{i,\ell}$。

表\ref{tab:q3features}给出候选增强模型使用的特征。所有状态、斜率和导数都在当前可见的$1{:}\ell$圈内重算；末端SOH/相对内阻取最后10圈中位数，导数取最后20圈中位数，近期趋势统一使用最近30圈。

\begin{table}[H]
\centering
\small
\rowcolors{2}{TableAlt}{white}
\begin{tabularx}{\textwidth}{lXl}
\toprule
\rowcolor{TableHeader}
分组 & 特征定义 & 作用 \\
\midrule
M1 & $\bar S_{\ell}$、SOH AUC、全局/近期Theil--Sen斜率及二者差值 & SOH水平与多尺度稳健趋势 \\
M1D & M1加最后20圈SG一阶导数中位数 & 局部瞬时退化速度 \\
M1D2 & M1D加最后20圈SG二阶导数中位数 & 局部曲率候选 \\
M2 & M1D加相对内阻末端水平、全局/近期斜率及差值；相对内阻以最初10圈中位数归一化 & 阻抗状态与增长趋势 \\
M2D & M2加相对内阻SG一阶导数 & 阻抗局部导数候选 \\
M3 & M2加温度、充电时间的中位状态、全局/近期斜率及差值 & 热与充电过程信息 \\
M4 & M3加九种策略哑变量 & 检验策略标签的边际信息 \\
\bottomrule
\end{tabularx}
\refstepcounter{table}\label{tab:q3features}
\end{table}

增强模型先计算每块训练电池相对TS基线的50维未来残差
\begin{equation}
\mathbf r_i=
\left(S_{i,\ell+1}-\widehat S_{i,\ell+1}^{\TS},\ldots,
S_{i,\ell+50}-\widehat S_{i,\ell+50}^{\TS}\right)^{\!\top}.
\end{equation}
将训练折的残差向量减去均值$\bar{\mathbf r}$并按行堆叠成矩阵$\mathbf R$，其协方差矩阵为
\begin{equation}
\mathbf C_r=\frac{1}{M-1}\mathbf R^{\top}\mathbf R.
\end{equation}
对$\mathbf C_r$作特征分解$\mathbf C_r\boldsymbol\phi_k=\lambda_k\boldsymbol\phi_k$，取最小的$K$使
\begin{equation}
\frac{\sum_{k=1}^{K}\lambda_k}{\sum_{k=1}^{50}\lambda_k}\ge0.95.
\end{equation}
于是残差和PCA得分分别为
\begin{equation}
\mathbf r_i\approx\bar{\mathbf r}+\sum_{k=1}^{K}z_{ik}\boldsymbol\phi_k,
\qquad
z_{ik}=(\mathbf r_i-\bar{\mathbf r})^{\top}\boldsymbol\phi_k.
\end{equation}

用Ridge回归\cite{hoerl1970}从标准化早期特征矩阵$\mathbf X$预测得分矩阵$\mathbf Z$：
\begin{equation}
J(\mathbf B)=\|\mathbf Z-\mathbf X\mathbf B\|_F^2
+\lambda\|\mathbf B\|_F^2.
\end{equation}
对$\mathbf B$求导并令其为0：
\begin{align}
\frac{\partial J}{\partial\mathbf B}
&=-2\mathbf X^{\top}(\mathbf Z-\mathbf X\mathbf B)+2\lambda\mathbf B=0,\\
(\mathbf X^{\top}\mathbf X+\lambda\mathbf I)\mathbf B
&=\mathbf X^{\top}\mathbf Z,\\
\widehat{\mathbf B}
&=(\mathbf X^{\top}\mathbf X+\lambda\mathbf I)^{-1}\mathbf X^{\top}\mathbf Z.
\end{align}
对新电池有$\widehat{\mathbf z}_i=\mathbf x_i^{\top}\widehat{\mathbf B}$、$\widehat{\mathbf r}_i=\bar{\mathbf r}+\sum_{k=1}^{K}\widehat z_{ik}\boldsymbol\phi_k$，最终预测为TS基线加残差修正。PCA、标准化、$K$和$\lambda$均只在每个LOBO训练折内确定。

\subsection{模型选择、消融与信息贡献}

在$\ell=150$的直接基线比较中，Persistence、OLS、TS和SG导数基线的RMSE依次为0.002233、0.0005343、0.0005289和0.0006730。TS只比OLS低约1.02\%，绝对差异很小。增强模型中M2的RMSE最低（0.000583），仍未超过TS；M3加入温度和充电时间后RMSE增至0.001882，M4加入策略标签后回落至0.000598，说明这些特征可解释部分残差，却没有形成稳定的样本外净收益。图\ref{fig:ablation}改用雷达图汇总七组候选：半径越小表示RMSE越低，珊瑚色虚线环为TS基线。七个候选点均位于基线环外，其中M2最接近基线，直观说明增加特征后并未获得可重复的样本外精度优势。

\begin{figure}[H]
\centering
\includegraphics[width=0.82\textwidth]{fig9_ablation.pdf}
\refstepcounter{figure}\label{fig:ablation}
\end{figure}

对任一内层候选$m$，以电池为独立单位计算$e_{im}=\mathrm{MSE}_{im}$，并定义
\begin{equation}
\bar e_m=\frac1M\sum_{i=1}^{M}e_{im},
\qquad
\mathrm{SE}_m=\frac{\mathrm{sd}(e_{1m},\ldots,e_{Mm})}{\sqrt M}.
\end{equation}
若$m^*=\arg\min_m\bar e_m$，则保留$\bar e_m\le\bar e_{m^*}+\mathrm{SE}_{m^*}$的近优集合，并选复杂度最低者。OLS与TS同为单斜率趋势基线；若二者同时进入近优集合，按运行前已写明的“TS为默认稳健基线”职责优先TS，且该次序不得根据外层误差改动。结果40个外层折均选择TS\_ONLY。旧的OLS优先规则同样40/40选择OLS，RMSE为0.0005343；正式TS规则为0.0005289。两者的敏感性对照说明“40/40选择TS”主要反映固定裁决规则，不能表述为TS逐折显著击败OLS。

\begin{table}[H]
\centering
\rowcolors{2}{TableAlt}{white}
\begin{tabular}{lr}
\toprule
\rowcolor{TableHeader}
指标 & 数值 \\
\midrule
外层模型选择 & TS\_ONLY，40/40折 \\
总体MAE & 0.000343 \\
总体RMSE & 0.000529 \\
总体MAPE/相对吻合度 & 0.03462\%/99.96538\% \\
总体$R^2$ & 0.9970 \\
第200圈绝对误差中位数/90\%分位数 & 0.000395/0.001056 \\
电池级RMSE中位数/90\%分位数/最大值 & 0.000277/0.000638/0.001841 \\
\bottomrule
\end{tabular}
\refstepcounter{table}\label{tab:q3nested}
\end{table}

因此，充电策略信息和IR在开发阶段可解释部分未来残差，但在严格模型选择后没有形成稳定的增量预测收益；温度和充电时间特征在当前样本规模下反而增加方差。对本题第151--200圈预测，最可靠的信息是电池自身最近30圈的实际SOH退化趋势。

不同可见长度下的固定Theil--Sen基线结果见表\ref{tab:length}和图\ref{fig:length}。从50圈增加到150圈，RMSE由0.00406降至0.000529，说明第100圈后的近期趋势对随后50圈预测尤为重要。这里$\ell=150$由题设可见数据长度确定，并非看完五组误差后择优，因此不产生“挑最优早期长度”的额外选择偏差。

\begin{table}[H]
\centering
\rowcolors{2}{TableAlt}{white}
\begin{tabular}{rrr}
\toprule
\rowcolor{TableHeader}
$\ell$/圈 & MAE & RMSE \\
\midrule
50 & 0.003339 & 0.004056 \\
75 & 0.001515 & 0.002315 \\
100 & 0.001062 & 0.001578 \\
125 & 0.001023 & 0.001229 \\
150 & \textbf{0.000343} & \textbf{0.000529} \\
\bottomrule
\end{tabular}
\refstepcounter{table}\label{tab:length}
\end{table}

\begin{figure}[H]
\centering
\includegraphics[width=0.92\textwidth]{fig8_length_accuracy.pdf}
\refstepcounter{figure}\label{fig:length}
\end{figure}

\subsection{测试电池预测与EOL外推}

将最终趋势模型用于9块测试电池，得到第151--200圈SOH，如图\ref{fig:testpred}。各预测曲线在第150圈边界连续，且对快速退化的16号电池给出明显更陡的下降趋势。

\begin{figure}[H]
\centering
\includegraphics[width=0.98\textwidth]{fig10_test_predictions.pdf}
\refstepcounter{figure}\label{fig:testpred}
\end{figure}

将每块测试电池的前150圈观测与未来50圈预测合并，并严格继承问题一的$n_0=31$中心化单调二次结构：
\begin{equation}
S_i(n)=A_i-B_i(n-31)-C_i(n-31)^2,qquad B_i,C_i\ge0.
\end{equation}
令$x_i=\widehat L_i^{\mathrm{ind}}-31$，由
\begin{align}
A_i-B_ix_i-C_ix_i^2&=0.8,\\
C_ix_i^2+B_ix_i-(A_i-0.8)&=0
\end{align}
得到个体交点
\begin{equation}
\widehat L_i^{\mathrm{ind}}=31+
\frac{-B_i+\sqrt{B_i^2+4C_i(A_i-0.8)}}{2C_i},\qquad C_i>0.
\end{equation}

为压制少数电池的曲率外推漂移，只在Q3九块测试电池上引入同策略部分池化。设$\widehat L_{p(i)}^{\mathrm{peer}}$为训练折内同策略电池的几何中心，定义
\begin{align}
\log \widehat L_i^{\mathrm{final}}
&=w\log \widehat L_i^{\mathrm{ind}}
+(1-w)\log \widehat L_{p(i)}^{\mathrm{peer}},\\
\widehat L_i^{\mathrm{final}}
&=(\widehat L_i^{\mathrm{ind}})^w
\cdot(\widehat L_{p(i)}^{\mathrm{peer}})^{1-w}.
\end{align}
候选$w$在每个外层训练折内选择，40/40折均取$w=0.75$，故正式点估计为75\%个体信息与25\%同策略训练信息的几何收缩。该步骤不回填Q1的49块统一基准寿命。

为传播151--200圈短期预测误差，按Bootstrap思想\cite{efron1979}从40块训练电池的整段未来残差中以“电池块”有放回抽样2000次，并对每次交点施加相同的训练同策略收缩；这样保留同一电池50圈残差的时间相关性。结果见表\ref{tab:testlife}。

\begin{table}[H]
\centering
\small
\begin{threeparttable}
\rowcolors{2}{TableAlt}{white}
\begin{tabular}{clrr}
\toprule
\rowcolor{TableHeader}
电池 & 快充策略 & 点估计/圈 & Bootstrap中位数及95\%区间/圈 \\
\midrule
2  & 3.6C--80\%--3.6C & 1651 & 1541 $[1115,2792]$ \\
5  & 80\%--3.6C & 1086 & 1056 $[907,1225]$ \\
9  & 4.8C--80\%--4.8C & 766 & 756 $[675,807]$ \\
10 & 5.0C--67\%--4.0C（新结构） & 1742 & 1616 $[1115,3492]$ \\
11 & 5.3C--54\%--4.0C（新结构） & 2754 & 2238 $[1253,5868]$ \\
14 & 5.6C--19\%--4.6C（新结构） & 1975 & 1772 $[1140,6478]$ \\
16 & 3.7C--31\%--5.9C（新结构） & 924 & 900 $[750,1049]$ \\
24 & 5.6C--36\%--4.3C（新结构） & 1798 & 1670 $[1186,3347]$ \\
25 & 4.8C--80\%--4.8C（新结构） & 1805 & 1671 $[1143,3915]$ \\
\bottomrule
\end{tabular}
\begin{tablenotes}
\footnotesize
\item 注：区间只反映训练电池第151--200圈预测残差传播到中心化二次交点并经固定权重收缩后的波动，未覆盖长期退化结构本身的不确定性；不能解释为真实EOL置信区间。
\end{tablenotes}
\end{threeparttable}
\refstepcounter{table}\label{tab:testlife}
\end{table}

9号和16号电池的点估计较低，11号最高但区间仍宽，体现缓慢趋势与二次曲率经远距离外推后的放大效应。这里不使用任何外部真实寿命映射校正。

在Q1的直接截断审计中，单电池中心化二次EOL的150--200圈更新中位数/均值为11.77\%/30.66\%，Spearman为0.847；同一外层框架加入25\%同策略收缩后，中位数/均值降至6.72\%/17.99\%。图\ref{fig:pooling}用逐电池配对线和两侧小提琴密度显示，多数电池的截断更新误差在部分池化后向较低区域移动，且极端长尾明显收缩，表明稳健化并非由少数电池单点改善驱动。图\ref{fig:eolaudit}进一步比较“Q3预测到200圈”与“真实观察到200圈”后的同结构外推，中位相对差异为17.23\%、Spearman仅0.654。三项比较都只是内部稳定性审计，二者均不是真实EOL标签；它们明确说明99.965\%的短期SOH吻合度不能外推成99\%量级的寿命准确率。

\begin{figure}[H]
\centering
\includegraphics[width=0.90\textwidth]{fig12_eol_structure_audit.pdf}
\refstepcounter{figure}\label{fig:eolaudit}
\end{figure}

\begin{figure}[H]
\centering
\includegraphics[width=0.90\textwidth]{fig12b_partial_pooling.pdf}
\refstepcounter{figure}\label{fig:pooling}
\end{figure}

另定义近期斜率变化$\Delta k_i=k_i^{151:200}-k_i^{121:150}$。在直接延续、全体均值、同策略均值和Ridge四个辅助头中，全体均值与同策略均值同处one-SE近优集合，故按简约原则取
\begin{equation}
\widehat{\Delta k}=\frac1{40}\sum_{i=1}^{40}\Delta k_i,qquad
\widehat k_i^{151:200}=k_i^{121:150}+\widehat{\Delta k}.
\end{equation}
其LOBO未来斜率RMSE由直接延续的$1.37245\times10^{-5}$降至$1.02947\times10^{-5}$，改善24.99\%。该头只说明近期退化存在总体加速修正，不参与EOL结构门控。

\subsection{长期三均值辅助压力测试}

题目还给出每块电池完整寿命区间的平均内阻、平均温度和平均充电时间。为避免长期信息泄漏，先完全冻结SOH预测和EOL，再只用40块训练电池前150圈逐循环数据拟合状态关系。对任一通道$X\in\{\mathrm{IR},T,\mathrm{ChargeTime}\}$，完整交互模型
\begin{equation}
X_{i,n}=\alpha_X+\beta_XS_{i,n}+\gamma_{X,P_i}
+\delta_{X,P_i}S_{i,n}+\varepsilon_{i,n}
\end{equation}
只用于检验Policy是否改变状态水平和SOH响应斜率，不承担正式轨迹外推。因为测试电池第150圈状态已观测，非平凡轨迹优先采用锚定形式
\begin{equation}
\widehat X_{i,n}=X_{i,150}^{\mathrm{rob}}
+(\beta_X+\delta_{X,P_i})
\left(\widehat S_{i,n}-S_{i,150}^{\mathrm{rob}}\right),\qquad n>150,
\end{equation}
其中全局锚定令$\delta_{X,P_i}=0$；Persistence令$\widehat X_{i,n}=X_{i,150}^{\mathrm{rob}}$。每次留出一块电池，先以真实151--200圈SOH检验状态关系，再以Q3预测SOH作端到端检查。

IR和温度的最优非平凡候选为全局锚定，充电时间为Policy锚定；相对Persistence的电池级MSE改善分别为5.88\%、6.88\%和2.39\%。但三者均落在最优误差的一标准误范围内，所以正式关系全部回到Persistence。这个结果表示复杂关系没有稳定泛化收益，不是模型失败，也没有理由继续叠加神经网络或更多交互。

令$\widetilde L_i=\max\{200,\operatorname{round}(\widehat L_i)\}$，把已观测1--150圈与冻结关系生成的151--$\widetilde L_i$圈拼接，完整寿命均值为
\begin{equation}
\widehat{\overline X}_i
=\frac{1}{\widetilde L_i}\sum_{n=1}^{\widetilde L_i}\widehat X_{i,n}.
\end{equation}
表\ref{tab:lifemeans}与图\ref{fig:lifemeans}给出与题目汇总真值的压力测试。

\begin{table}[H]
\centering
\rowcolors{2}{TableAlt}{white}
\begin{tabular}{lrrrr}
\toprule
\rowcolor{TableHeader}
状态量 & MAE & 相对MAE & 平均偏差 & Spearman $\rho$ \\
\midrule
平均内阻 & 0.000206 & 1.261\% & $-0.000206$ & 1.000 \\
平均温度 & $0.605\,^{\circ}\mathrm C$ & 1.751\% & $-0.605\,^{\circ}\mathrm C$ & 0.983 \\
平均充电时间 & $0.535\,\mathrm{min}$ & 4.347\% & $-0.523\,\mathrm{min}$ & 0.583 \\
\bottomrule
\end{tabular}
\refstepcounter{table}\label{tab:lifemeans}
\end{table}

\begin{figure}[H]
\centering
\includegraphics[width=0.99\textwidth]{fig11_life_summary_validation.pdf}
\refstepcounter{figure}\label{fig:lifemeans}
\end{figure}

三通道相对吻合度$100\%-\mathrm{MAPE}$分别为98.739\%、98.249\%和95.653\%，但这些数字衡量的是完整寿命\emph{状态均值}，不是EOL圈数。Persistence关系使状态均值对寿命长度不敏感；即使IR和温度排序相关很高，也不能据此校准$\widehat L_i$。结合图\ref{fig:eolaudit}的17.23\%结构差异与仅0.654的排序相关，正确结论是：151--200圈短期SOH预测获得了直接标签支持，长期EOL仍缺少真实标签和可靠跨通道验证。

\section{四问逐问作答核对}

表\ref{tab:coverage}按题面逐项核对输出，确保每个小问均有方法、数值结果和证据边界。

\begin{table}[H]
\centering
\footnotesize
\rowcolors{2}{TableAlt}{white}
\begin{tabularx}{\textwidth}{cclX}
\toprule
\rowcolor{TableHeader}
问题 & 小问 & 对应位置 & 明确回答 \\
\midrule
Q1 & 1 & 表\ref{tab:q1battery}、图\ref{fig:policysoh} & 逐电池列出策略、充电时间、SOH状态/斜率及寿命 \\
Q1 & 2 & 表\ref{tab:q1policy}、图\ref{fig:lifedist} & 给出九策略寿命中位数、IQR及电池级分布 \\
Q1 & 3 & 表\ref{tab:longshortdetail} & 按EOL口径正式给出典型长、短寿命策略，早期指标均作一致性审计 \\
Q1 & 4 & 表\ref{tab:longshortdetail}、图\ref{fig:longshort} & 量化SOH、稳定速率、AUC、时间和温度，并说明跨数据集不作参数因果解释 \\
\midrule
Q2 & 1 & 策略总体差异小节 & 稳定速率$p=0.000200$；EOL $p=0.01522$；分别报告Dunn--Holm结果 \\
Q2 & 2 & 参数回归小节 & 报告$C_1,q,C_2$相关、理论时间、VIF及$A,D$重参数化 \\
Q2 & 3 & 表\ref{tab:q2reg} & 主响应精确$p=13/720$且LOPO 6/6同向；EOL辅助证据有限 \\
Q2 & 4 & 图\ref{fig:socsens} & 40\%--60\%分界下中高SOC正系数持续更大 \\
Q2 & 5 & 表\ref{tab:mechanism} & 总体倍率与热/阻抗线索相容；SOC位置通道未被识别并列明局限 \\
\midrule
Q3 & 1 & 表\ref{tab:q3features} & 明确定义SOH、IR、温度、时间及策略特征与窗口 \\
Q3 & 2 & 趋势基线与增强模型小节 & 建立近期TS趋势与Trend--PCA--Ridge候选模型 \\
Q3 & 3 & 图\ref{fig:testpred}、表\ref{tab:testlife} & 输出9块电池151--200圈轨迹及部分池化中心化二次EOL \\
Q3 & 4 & 表\ref{tab:q3nested}、图\ref{fig:eolaudit} & 短期RMSE为0.000529；远期EOL仍只作截断稳定性审计 \\
Q3 & 5 & 图\ref{fig:ablation}、图\ref{fig:pooling}、图\ref{fig:lifemeans} & 给出特征消融、部分池化及三状态通道的实际增量/负结果 \\
\midrule
Q4 & 1 & 时间模型小节、表\ref{tab:q4pareto} & 由恒流公式推导理论时间，并用九策略观测时间校验字段口径 \\
Q4 & 2 & 表\ref{tab:q4pareto}、图\ref{fig:q4pareto} & 以直接观测稳定退化速率估计九种已有策略的损伤 \\
Q4 & 3 & 局部鲁棒优化小节 & 建立时间预算、双重凸包、近邻约束和Bootstrap分位数目标 \\
Q4 & 4 & 表\ref{tab:q4recommend} & 推荐N2，并与正式长寿命N6及短寿命S3比较 \\
Q4 & 5 & Q3一致性与风险小节 & 给出适用域、配对Bootstrap、短期未来审计及新策略实验要求 \\
\bottomrule
\end{tabularx}
\refstepcounter{table}\label{tab:coverage}
\end{table}

\section{问题四：充电策略优化}

问题四先在九种真实实验策略上建立观测时间--稳定退化的经验帕累托前沿，再把连续搜索严格限制在同类策略的局部数据支持域内，并用独立自助重采样检验新候选是否真正优于最佳已有策略。图\ref{fig:workflowq4}概括了从真实策略优选到局部鲁棒优化、再到最终推荐裁决的完整路径。

\begin{figure}[H]
\centering
\includegraphics[width=0.98\textwidth]{fig_workflow_q4.pdf}
\refstepcounter{figure}\label{fig:workflowq4}
\end{figure}

\subsection{优化目标、时间模型与主退化指标}

令策略向量为$\boldsymbol\theta=(C_1,q,C_2)$。若恒流倍率为$C$、充入SOC比例为$\Delta s$，由“时间=容量/电流”得到$60\Delta s/C$分钟，故两阶段相加为
\begin{align}
T_1(\boldsymbol\theta)&=60\frac{q}{C_1},\\
T_2(\boldsymbol\theta)&=60\frac{0.8-q}{C_2},\\
T_{0\rightarrow80}^{\mathrm{theo}}(\boldsymbol\theta)
&=T_1+T_2
=60\left(\frac{q}{C_1}+\frac{0.8-q}{C_2}\right).
\end{align}
逐循环\texttt{chargetime}按前述口径记作$T_{i}^{\mathrm{obs}}$。已有策略$g$的观测时间与退化分别汇总为
\begin{equation}
T_g^{\mathrm{obs}}=\operatorname{median}_{i\in g}T_i^{\mathrm{obs}},
\qquad
R_g=\operatorname{median}_{i\in g}r_i^{\mathrm{stable}}.
\end{equation}
概念上的双目标为
\begin{equation}
\min_{\boldsymbol\theta}\left(
T_{0\rightarrow80}(\boldsymbol\theta),
r^{\mathrm{stable}}(\boldsymbol\theta)
\right).
\end{equation}
但Q1/Q3的80\% EOL没有真实标签，且“预测至200圈/真实观察至200圈”后的同结构外推中位差仍为17.23\%、排序相关仅0.654，故问题四不把EOL点估计作为连续优化主目标，而以直接观测的$r^{\mathrm{stable}}$表示寿命损伤。已有策略比较使用$T_g^{\mathrm{obs}}$；新参数没有实测时间，只把物理公式作为设计约束，并单列其理论--实测偏差风险。

\subsection{九种已有策略的经验Pareto优选}

对两个均需最小化的指标，若策略$g_1$满足
\begin{equation}
T_{g_1}^{\mathrm{obs}}\le T_{g_2}^{\mathrm{obs}},\qquad
R_{g_1}\le R_{g_2},
\end{equation}
且至少一个不等号严格成立，则称$g_1$支配$g_2$。不被任何其他策略支配的集合为
\begin{equation}
\mathcal P_{\mathrm{obs}}
=\left\{g:\nexists h\ne g, h\prec g\right\}.
\end{equation}
表\ref{tab:q4pareto}使用全部49块电池前150圈数据。为评估小样本下前沿身份的波动，另在每个策略内部有放回重采样5000次，并统计每次进入Pareto集的频率；该频率是稳定性描述，不是显著性$p$值。

\begin{table}[H]
\centering
\footnotesize
\rowcolors{2}{TableAlt}{white}
\begin{tabular}{clrrrrc}
\toprule
\rowcolor{TableHeader}
代码 & 策略 & $N$ & $T_g^{\mathrm{obs}}$/min & $10^4R_g$ & 非支配频率 & 经验Pareto \\
\midrule
S1 & 3.6C--80\%--3.6C & 3 & 13.342 & 0.197 & 36.64\% & 否 \\
S2 & 80\%--3.6C & 3 & 13.342 & 0.272 & 0.00\% & 否 \\
S3 & 4.8C--80\%--4.8C（原结构） & 3 & 10.054 & 0.849 & 25.58\% & 否 \\
N1 & 5.0C--67\%--4.0C（新结构） & 7 & 10.043 & 0.231 & 52.78\% & 否 \\
N2 & 5.3C--54\%--4.0C（新结构） & 8 & 10.043 & 0.217 & 74.66\% & 是 \\
N3 & 5.6C--19\%--4.6C（新结构） & 8 & 10.043 & 0.306 & 12.44\% & 否 \\
N4 & 3.7C--31\%--5.9C（新结构） & 3 & 10.096 & 1.925 & 0.00\% & 否 \\
N5 & 5.6C--36\%--4.3C（新结构） & 8 & 10.043 & 0.273 & 57.42\% & 是 \\
N6 & 4.8C--80\%--4.8C（新结构） & 6 & 11.038 & 0.181 & 69.10\% & 是 \\
\bottomrule
\end{tabular}
\refstepcounter{table}\label{tab:q4pareto}
\end{table}

\begin{figure}[H]
\centering
\includegraphics[width=0.98\textwidth]{fig13_q4_empirical_pareto.pdf}
\refstepcounter{figure}\label{fig:q4pareto}
\end{figure}

精确支配下，N2、N5与N6构成三点前沿。N5比N2仅快$0.000147$ min（约0.009 s），但观测退化速率高21.6\%，这一时间差不具工程分辨意义；N6的观测退化速率比N2低18.8\%，代价是约0.995 min更长的记录时间。因此，在强调约10分钟快充时N2是更合理的已有策略；若可接受约1 min额外时间并优先压低直接观测退化，N6是备选。

\subsection{六个近等时新结构策略的局部可行域}

连续优化只使用六个同时属于数据集3和新结构类别的参数位置。记其原始参数点为$\boldsymbol\theta_j$，$j=1,\ldots,6$，先定义原始参数凸包
\begin{equation}
\Omega_{\theta}=\operatorname{conv}\{\boldsymbol\theta_1,\ldots,\boldsymbol\theta_6\}.
\end{equation}
为避免“原始参数看似插值，但进入Q2回归后在$(A,D)$空间外推”，再令$\Phi(\boldsymbol\theta)=(A(\boldsymbol\theta),D(\boldsymbol\theta))$并要求
\begin{equation}
\Phi(\boldsymbol\theta)\in
\Omega_{AD}=\operatorname{conv}\{\Phi(\boldsymbol\theta_1),\ldots,\Phi(\boldsymbol\theta_6)\}.
\end{equation}
对原始参数按六点均值$\boldsymbol\mu_\theta$和标准差$\boldsymbol s_\theta$逐维标准化：
\begin{equation}
\boldsymbol z(\boldsymbol\theta)
=\operatorname{diag}(\boldsymbol s_\theta)^{-1}
(\boldsymbol\theta-\boldsymbol\mu_\theta).
\end{equation}
候选到最近真实策略的距离及允许上限定义为
\begin{align}
d(\boldsymbol\theta)&=\min_j
\|\boldsymbol z(\boldsymbol\theta)-\boldsymbol z(\boldsymbol\theta_j)\|_2,\\
d_{\max}&=\max_j\min_{k\ne j}
\|\boldsymbol z(\boldsymbol\theta_j)-\boldsymbol z(\boldsymbol\theta_k)\|_2=3.326.
\end{align}
于是局部可行域为
\begin{equation}
\boxed{\Omega_{\mathrm{local}}=\left\{
\boldsymbol\theta\in\Omega_{\theta}:
\Phi(\boldsymbol\theta)\in\Omega_{AD},\ d(\boldsymbol\theta)\le d_{\max}
\right\}}.
\end{equation}
采用可实际执行的分辨率$C_1,C_2\in0.1\mathbb Z$、$q\in0.01\mathbb Z$，并取六个实验设计中最大的理论时间为预算：
\begin{equation}
\mathcal G=\left\{\boldsymbol\theta\in\Omega_{\mathrm{local}}:
T_{0\rightarrow80}^{\mathrm{theo}}(\boldsymbol\theta)\le10.0132\ \mathrm{min}
\right\}.
\end{equation}
双重凸包、近邻距离和时间约束后共保留597个网格点。该集合只允许局部插值，不能扩展为任意倍率组合。

\subsection{Bootstrap分位数目标与独立稳定性审计}

问题二模型使用标准化$A^*,D^*$。在第$b$次Bootstrap中，对每个策略内部按电池有放回重采样，重新计算六个策略中位退化速率并拟合
\begin{equation}
\widehat r^{(b)}(\boldsymbol\theta)
=\widehat\beta_0^{(b)}+
\widehat\beta_A^{(b)}A^*(\boldsymbol\theta)+
\widehat\beta_D^{(b)}D^*(\boldsymbol\theta).
\end{equation}
不用均值最漂亮的点，而以较不利侧90\%分位数为选择目标：
\begin{equation}
J_{0.90}(\boldsymbol\theta)
=Q_{0.90}\left(\widehat r^{(1)}(\boldsymbol\theta),\ldots,
\widehat r^{(5000)}(\boldsymbol\theta)\right),
\qquad
\boldsymbol\theta^*=\arg\min_{\boldsymbol\theta\in\mathcal G}J_{0.90}(\boldsymbol\theta).
\end{equation}
为避免在同一批重采样上选择又评价，另生成独立5000次Bootstrap作稳定性审计。令$\boldsymbol\theta_E$为选择阶段最佳已有策略N2，定义配对差
\begin{equation}
\Delta^{(b)}=\widehat r^{(b)}(\boldsymbol\theta^*)-
\widehat r^{(b)}(\boldsymbol\theta_E).
\end{equation}
只有同时满足$Q_{0.90}(\Delta)<0$，且中位改善超过Q2策略级LOPO MAE，才把新点升级为正式推荐；这一裁决在查看验证结果前固定。

局部搜索得到$\boldsymbol\theta^*=(5.1\mathrm C,45\%,4.5\mathrm C)$，其理论时间为9.9608 min，$A=4.8375$、$D=-0.540$、最近标准化距离为0.930。图\ref{fig:q4local}左侧以三维参数空间展示597个可行网格及真实策略位置，并显式标出坐标原点$O(0,0,0)$；原点仅用于空间参照，不代表最优方向。右侧再映射到$(A,D)$空间并用等值热区表示退化预测的90\%分位数：由于问题二中$A,D$系数均为正，在当前局部模型内越靠近左下方（$A,D$更小）预测退化越慢，且颜色越浅表示90\%分位数越低。候选点处于已有设计邻域内，而非参数长方体边界的任意拼接。

\begin{figure}[H]
\centering
\includegraphics[width=0.98\textwidth]{fig14_q4_local_robust_search.pdf}
\refstepcounter{figure}\label{fig:q4local}
\end{figure}

独立验证中，N2与局部候选的预测分位数见表\ref{tab:q4robust}。候选的边际$Q_{0.90}$较低，但配对比较显示，其中位改善为$-2.43\times10^{-7}$，即候选中位预测反而略高；$P(\Delta<0)=47.1\%$，$Q_{0.90}(\Delta)=3.24\times10^{-6}>0$。这些差异也没有超过LOPO MAE $7.91\times10^{-6}$，故不满足新策略推荐门槛。换言之，5000次独立配对Bootstrap的差值分布跨越零点，左侧“候选更优”的概率仅47.1\%，而90\%分位数仍落在零点右侧，因此不升级为正式推荐。

\begin{table}[H]
\centering
\small
\rowcolors{2}{TableAlt}{white}
\begin{tabular}{lrrrrrrrr}
\toprule
\rowcolor{TableHeader}
方案 & $C_1$ & $q$/\% & $C_2$ & $T^{\mathrm{theo}}$/min & $A$ & $D$ & $10^4Q_{0.50}$ & $10^4Q_{0.90}$ \\
\midrule
N2（已有） & 5.3 & 54 & 4.0 & 10.0132 & 4.8775 & $-1.1267$ & 0.204 & 0.239 \\
C$^*$（待验证） & 5.1 & 45 & 4.5 & 9.9608 & 4.8375 & $-0.5400$ & 0.206 & 0.227 \\
\bottomrule
\end{tabular}
\refstepcounter{table}\label{tab:q4robust}
\end{table}

\subsection{最终推荐及与典型长短寿命策略比较}

综合经验Pareto、独立Bootstrap和LOPO可分辨尺度，\textbf{最终推荐已有实验策略N2：5.3C--54\%--4.0C新结构}。C$^*$仅作为后续实验候选，不应在缺少真实循环数据时宣称优于N2。表\ref{tab:q4recommend}还对照了问题一正式典型长寿命N6和典型短寿命S3。

\begin{table}[H]
\centering
\footnotesize
\rowcolors{2}{TableAlt}{white}
\begin{tabular}{clrrrrr}
\toprule
\rowcolor{TableHeader}
代码 & 角色 & $T^{\mathrm{obs}}$/min & $T^{\mathrm{theo}}$/min & $10^4R_g$ & EOL/圈 & $10^4Q_{0.90}$ \\
\midrule
N2 & 最终推荐已有策略 & 10.043 & 10.013 & 0.217 & 1593 & 0.239 \\
N6 & 正式典型长寿命 & 11.038 & 10.000 & 0.181 & 1775 & 0.242 \\
S3 & 典型短寿命（原结构） & 10.054 & 10.000 & 0.849 & 640 & -- \\
C$^*$ & 局部待验证候选 & -- & 9.961 & -- & -- & 0.227 \\
\bottomrule
\end{tabular}
\refstepcounter{table}\label{tab:q4recommend}
\end{table}

N6的观测退化速率比N2低约16.5\%，且基准EOL更高，但观测充电时间约多0.995 min；因此在强调约10分钟快充时N2更合适，若允许约1 min额外时间并优先压低退化，N6是备选。N2相对S3的观测退化速率低约74.5\%，基准EOL约为其2.49倍。S3虽与N6具有相同数值参数，却属于不同数据集和结构类别，因此数据集3的$(A,D)$回归没有被错误地移用到S3，表中不报告其鲁棒预测分位数。

\subsection{问题三短期一致性与适用风险}

对40块完整电池，取外层OOF在第200圈的预测，并计算第151--200圈观测与预测平均退化速率。九策略的$r_g^{\mathrm{stable}}$与真实未来退化速率的Spearman相关为0.750（$p=0.0199$），与OOF预测未来退化速率的相关为0.833（$p=0.00527$）；样本仅九个策略，故这里只作方向一致性审计。N2在真实未来退化速率中排第2，N6排第1，均未与问题四的低退化判断发生明显冲突。

全新策略C$^*$没有前150圈SOH，不能直接输入问题三。合理部署流程是：先把C$^*$作为实验候选运行早期循环；获得至少150圈真实SOH后，再用问题三最近30圈Theil--Sen模型更新第151--200圈短期健康预测。最终结论受以下边界限制：
\begin{enumerate}[label=\arabic*)]
  \item 连续响应面只有六个同类参数位置，局部候选仍可能受未建模交互和批次效应影响；
  \item 理论时间适合约束设计，但4.8C--80\%--4.8C新结构存在1.038 min理论--实测偏差，新候选的真实时间必须实验确认；
  \item \texttt{NEWSTRUCTURE}只是数据类别标签，不能外推为所有新结构电芯的普遍规律；
  \item 推荐主要降低前150圈稳定退化速率，不等同真实80\% EOL已经得到优化；
  \item N2是当前数据内的主推荐，N6是允许约1 min额外时间时的直接观测低退化备选，C$^*$只适合进一步实验筛查。
\end{enumerate}

\section{模型评价与当前结论边界}

\subsection{优点}

\begin{enumerate}[label=\arabic*)]
  \item 预处理只修复有充分局部证据的孤立尖峰，避免把真实退化误删为异常；
  \item 明确隔离第1--150圈建模与第151--200圈验证，Q1与Q3均用外层LOBO评价完整选型流程，防止同一电池未来信息参与模型选择；
  \item Q1把稳定起点与线性、幂律、中心化二次三类结构共同开放选择，并用外层LOBO评价完整流程；
  \item Q3的同策略部分池化权重在40个外层折内均选择0.75，同时降低EOL更新差的中位数和均值，主要抑制尾部漂移；
  \item 参数效应以策略为统计单位，并限制在同一结构、同一数据集内，减少伪重复和批次混杂；
  \item 增强模型未形成超过一标准误的稳定收益时保留简单模型；Q3的40个外层折均选择TS\_ONLY，结果简洁、可解释且可复现；
  \item Q4先回答已有实验策略的经验Pareto，再以原始参数和$(A,D)$双重凸包限制连续搜索，并用独立Bootstrap与LOPO误差尺度阻止不稳定的新策略推荐。
\end{enumerate}

\subsection{局限与后续工作}

\begin{enumerate}[label=\arabic*)]
  \item 观测SOH大多远高于0.8，寿命圈数对极小斜率误差高度敏感，故绝对EOL只能称为基准外推；
  \item 九种策略跨三个数据集，且部分策略仅3块电池；跨数据集的长短排序可能同时包含批次和结构影响；
  \item 数据集3仅有六个不同参数位置、残差自由度为3。主响应精确置换虽达到$p=0.0181$且LOPO方向稳定，但不能据此作全参数空间的精确因果识别；EOL辅助回归虽有$p=0.0458$，Bootstrap区间跨0且预测不稳；
  \item Q3的寿命Bootstrap只传播未来50圈经验残差，没有覆盖退化函数的结构不确定性；因此Q4采用直接观测稳定退化速率而非EOL点估计作主优化目标。
  \item Q3的嵌套LOBO是题目附件内部验证，尚未在新的实验批次、温度条件或电芯类型上作外部验证；$R^2=0.9970$只支持151--200圈短期SOH预测，不代表真实EOL预测精度。
  \item 完整寿命三均值的95.65\%--98.74\%吻合度不等于EOL准确率；两种“预测至200圈/观察至200圈”同结构外推中位差仍为17.23\%、Spearman仅0.654，不能替代真实EOL标签。
  \item Q4局部候选的理论时间与退化均来自模型；虽然采用双重凸包和近邻约束，仍须以新的电池实验验证，不能直接用于实际充电控制。
\end{enumerate}

综上，前150圈能够区分明显的快、慢退化策略，并以约99.965\%的SOH相对吻合度预测随后50圈；在六个近等时新结构策略位置上，较高平均倍率和更偏向中高SOC的倍率暴露与更快稳定退化相关。九种真实策略的时间--退化Pareto前沿包含N2、N5和N6；约10分钟预算下最终推荐已有实验策略N2，即5.3C--54\%--4.0C新结构。局部候选C$^*$未通过独立配对Bootstrap与LOPO尺度检验，故只作为待实验方案。Q1/Q3的中心化二次EOL经部分池化后内部稳定性有所提高，但仍只是一致口径下的情景指标，长期结构没有真实标签支持；问题四的推荐降低的是早期稳定退化风险，不能等同真实寿命已经被精确优化。

\begin{thebibliography}{99}
\bibitem{severson2019}
Severson K A, Attia P M, Jin N, et al. Data-driven prediction of battery cycle life before capacity degradation[J]. Nature Energy, 2019, 4: 383--391. DOI: 10.1038/s41560-019-0356-8.

\bibitem{savitzky1964}
Savitzky A, Golay M J E. Smoothing and differentiation of data by simplified least squares procedures[J]. Analytical Chemistry, 1964, 36(8): 1627--1639. DOI: 10.1021/ac60214a047.

\bibitem{sen1968}
Sen P K. Estimates of the regression coefficient based on Kendall's tau[J]. Journal of the American Statistical Association, 1968, 63(324): 1379--1389. DOI: 10.1080/01621459.1968.10480934.

\bibitem{kruskal1952}
Kruskal W H, Wallis W A. Use of ranks in one-criterion variance analysis[J]. Journal of the American Statistical Association, 1952, 47(260): 583--621. DOI: 10.1080/01621459.1952.10483441.

\bibitem{dunn1964}
Dunn O J. Multiple comparisons using rank sums[J]. Technometrics, 1964, 6(3): 241--252. DOI: 10.1080/00401706.1964.10490181.

\bibitem{holm1979}
Holm S. A simple sequentially rejective multiple test procedure[J]. Scandinavian Journal of Statistics, 1979, 6(2): 65--70.

\bibitem{hoerl1970}
Hoerl A E, Kennard R W. Ridge regression: biased estimation for nonorthogonal problems[J]. Technometrics, 1970, 12(1): 55--67. DOI: 10.1080/00401706.1970.10488634.

\bibitem{efron1979}
Efron B. Bootstrap methods: another look at the jackknife[J]. The Annals of Statistics, 1979, 7(1): 1--26. DOI: 10.1214/aos/1176344552.

\bibitem{ansean2016}
Anse\'an D, Dubarry M, Devie A, et al. Fast charging technique for high power LiFePO$_4$ batteries: A mechanistic analysis of aging[J]. Journal of Power Sources, 2016, 321: 201--209. DOI: 10.1016/j.jpowsour.2016.04.140.

\bibitem{gao2017}
Gao Y, Jiang J, Zhang C, et al. Lithium-ion battery aging mechanisms and life model under different charging stresses[J]. Journal of Power Sources, 2017, 356: 103--114. DOI: 10.1016/j.jpowsour.2017.04.084.
\end{thebibliography}

\newpage
\section*{附录：复现说明}

\subsection*{A. 逐电池关键结果}

\begin{longtable}{rrlrrrr}
\caption{49块电池的策略、充电时间、SOH与基准估计寿命}\label{tab:q1battery}\\
\toprule
电池 & 数据集 & 快充策略 & $T_{1:150}$/min & $S_{150}$ & $10^4k^{\TS}$ & $L$/圈 \\
\midrule
\endfirsthead
\multicolumn{7}{c}{表\thetable\ （续）}\\
\toprule
电池 & 数据集 & 快充策略 & $T_{1:150}$/min & $S_{150}$ & $10^4k^{\TS}$ & $L$/圈 \\
\midrule
\endhead
\midrule\multicolumn{7}{r}{续下页}\\\endfoot
\bottomrule\endlastfoot
1 & 1 & 3.6C--80\%--3.6C & 13.358 & 1.00303 & -0.197 & 1070 \\
2 & 1 & 3.6C--80\%--3.6C & 13.342 & 0.99533 & -0.200 & 1709 \\
3 & 1 & 3.6C--80\%--3.6C & 13.342 & 0.99624 & -0.176 & 1389 \\
4 & 2 & 80\%--3.6C & 13.374 & 0.97758 & -1.586 & 629 \\
5 & 2 & 80\%--3.6C & 13.342 & 0.99638 & -0.272 & 1114 \\
6 & 2 & 80\%--3.6C & 13.342 & 0.99634 & -0.268 & 1548 \\
7 & 2 & 4.8C--80\%--4.8C & 10.054 & 0.99511 & -0.406 & 808 \\
8 & 2 & 4.8C--80\%--4.8C & 10.093 & 0.99021 & -0.849 & 612 \\
9 & 2 & 4.8C--80\%--4.8C & 10.021 & 0.98955 & -0.938 & 640 \\
10 & 3 & 5.0C--67\%--4.0C（新结构） & 10.042 & 0.99721 & -0.231 & 2843 \\
11 & 3 & 5.3C--54\%--4.0C（新结构） & 10.041 & 0.99682 & -0.267 & 7638 \\
12 & 3 & 5.6C--19\%--4.6C（新结构） & 10.044 & 0.99641 & -0.282 & 7144 \\
13 & 3 & 5.6C--36\%--4.3C（新结构） & 10.041 & 0.99696 & -0.255 & 7894 \\
14 & 3 & 5.6C--19\%--4.6C（新结构） & 10.046 & 0.99765 & -0.204 & 1962 \\
15 & 3 & 5.6C--36\%--4.3C（新结构） & 10.042 & 0.99616 & -0.305 & 1434 \\
16 & 3 & 3.7C--31\%--5.9C（新结构） & 10.096 & 0.97277 & -1.925 & 856 \\
17 & 3 & 4.8C--80\%--4.8C（新结构） & 11.039 & 0.99833 & -0.127 & 4605 \\
18 & 3 & 5.0C--67\%--4.0C（新结构） & 10.043 & 0.99710 & -0.278 & 1340 \\
19 & 3 & 5.3C--54\%--4.0C（新结构） & 10.042 & 0.99845 & -0.135 & 2127 \\
20 & 3 & 4.8C--80\%--4.8C（新结构） & 11.038 & 0.99771 & -0.216 & 1659 \\
21 & 3 & 5.6C--19\%--4.6C（新结构） & 10.043 & 0.99582 & -0.368 & 1234 \\
22 & 3 & 5.6C--36\%--4.3C（新结构） & 10.043 & 0.99797 & -0.176 & 4155 \\
23 & 3 & 5.6C--19\%--4.6C（新结构） & 10.042 & 0.99604 & -0.369 & 1261 \\
24 & 3 & 5.6C--36\%--4.3C（新结构） & 10.042 & 0.99767 & -0.213 & 1305 \\
25 & 3 & 4.8C--80\%--4.8C（新结构） & 11.040 & 0.99730 & -0.230 & 8754 \\
26 & 3 & 5.3C--54\%--4.0C（新结构） & 10.042 & 0.99601 & -0.222 & 1313 \\
27 & 3 & 5.6C--19\%--4.6C（新结构） & 10.043 & 0.99574 & -0.348 & 5816 \\
28 & 3 & 5.6C--36\%--4.3C（新结构） & 10.043 & 0.99572 & -0.357 & 5648 \\
29 & 3 & 5.0C--67\%--4.0C（新结构） & 10.043 & 0.99731 & -0.250 & 1504 \\
30 & 3 & 3.7C--31\%--5.9C（新结构） & 10.091 & 0.97922 & -1.549 & 1307 \\
31 & 3 & 5.0C--67\%--4.0C（新结构） & 10.042 & 0.99760 & -0.212 & 1122 \\
32 & 3 & 5.3C--54\%--4.0C（新结构） & 10.043 & 0.99754 & -0.205 & 1133 \\
33 & 3 & 5.6C--19\%--4.6C（新结构） & 10.042 & 0.99737 & -0.230 & 1276 \\
34 & 3 & 5.6C--36\%--4.3C（新结构） & 10.045 & 0.99657 & -0.292 & 1096 \\
35 & 3 & 3.7C--31\%--5.9C（新结构） & 10.112 & 0.97188 & -2.167 & 636 \\
36 & 3 & 5.3C--54\%--4.0C（新结构） & 10.044 & 0.99753 & -0.211 & 1420 \\
37 & 3 & 5.0C--67\%--4.0C（新结构） & 10.043 & 0.99657 & -0.298 & 1579 \\
38 & 3 & 5.3C--54\%--4.0C（新结构） & 10.044 & 0.99579 & -0.357 & 1661 \\
39 & 3 & 5.6C--19\%--4.6C（新结构） & 10.044 & 0.99643 & -0.290 & 2319 \\
40 & 3 & 5.6C--36\%--4.3C（新结构） & 10.043 & 0.99540 & -0.374 & 1752 \\
41 & 3 & 4.8C--80\%--4.8C（新结构） & 11.038 & 0.95167 & -0.314 & 1295 \\
42 & 3 & 5.0C--67\%--4.0C（新结构） & 10.044 & 0.99755 & -0.200 & 1938 \\
43 & 3 & 5.3C--54\%--4.0C（新结构） & 10.044 & 0.99640 & -0.312 & 1526 \\
44 & 3 & 5.6C--19\%--4.6C（新结构） & 10.043 & 0.99654 & -0.323 & 1358 \\
45 & 3 & 5.6C--36\%--4.3C（新结构） & 10.043 & 0.99826 & -0.147 & 1302 \\
46 & 3 & 4.8C--80\%--4.8C（新结构） & 11.038 & 0.99836 & -0.146 & 1514 \\
47 & 3 & 5.0C--67\%--4.0C（新结构） & 10.044 & 0.99840 & -0.146 & 1321 \\
48 & 3 & 5.3C--54\%--4.0C（新结构） & 10.043 & 0.99829 & -0.162 & 1830 \\
49 & 3 & 4.8C--80\%--4.8C（新结构） & 11.038 & 0.99833 & -0.143 & 1891 \\
\end{longtable}


\subsection*{B. 论文与图形复现文件}

本次整理后的交付包保留主论文、编译类文件、绘图脚本、绘图所需结果数据及全部图形导出文件：
\begin{itemize}
  \item 主论文\texttt{.tex}与编译后的\texttt{.pdf}；
  \item \texttt{cumcmthesis.cls}：论文版式类文件；
  \item \texttt{figures\_generated\_v5/}：全部中文学术图，每幅图同时提供可编辑矢量PDF、SVG和高分辨率PNG；
  \item \texttt{plot\_all\_figures.py}：统一配色、中文字体、流程图、热图、小提琴图、密度图、三维参数域和帕累托图的绘制脚本；
  \item \texttt{data\_used/}：绘图脚本直接读取的正式结果CSV，保证图形可独立复现；
  \item \texttt{README\_图表与编译说明.txt}：XeLaTeX编译方式、绘图复现命令和文件说明。
\end{itemize}

\subsection*{C. 关键计算顺序}

\begin{enumerate}[label=\arabic*)]
  \item 读取两个题目CSV并检查电池--循环唯一性；
  \item 按电池局部Hampel/MAD标记异常，仅修复满足局部孤立性条件的容量、内阻和温度尖峰；
  \item 重算SOH并做SG(11,2)平滑；
  \item 联合比较五个稳定起点与三类低参数寿命结构，以151--200圈误差形成one-SE近优集合，再按无效EOL、截断更新、边界、复杂度和起点裁决，并用外层LOBO评价完整选型流程；
  \item 汇总九种策略，完成KW置换检验、Dunn--Holm比较和数据集3策略级参数回归；
  \item 在40块训练电池上完成外层LOBO、内层候选选择、一标准误简约决策、开发阶段特征消融和早期长度敏感性，再冻结模型预测9块测试电池；
  \item 用训练电池未来残差按电池块重采样，计算测试电池个体EOL，再以75\%个体、25\%同策略训练中心作对数尺度部分池化；
  \item 只用训练电池前150圈拟合三种状态关系并作电池级LOBO；在SOH/EOL冻结后才与三个完整寿命均值比较，且不据此反向校准寿命；
  \item 以九策略观测时间和稳定退化速率构造经验Pareto前沿；
  \item 在六个同类参数位置的原始参数凸包与$(A,D)$特征凸包交集中生成可实施网格，并施加近邻距离和10.0132 min理论时间预算；
  \item 用5000次Bootstrap选择$Q_{0.90}$最小候选，再用独立5000次Bootstrap、配对差和Q2 LOPO误差尺度决定是否升级为正式推荐；
  \item 运行Python绘图脚本生成全部中文学术图形，再由XeLaTeX编译论文。
\end{enumerate}

\end{document}
